% Praxis: symbolic reasoning between a language model and an enterprise decision.
% Build: pdflatex main && bibtex main && pdflatex main && pdflatex main
% Every number in this paper comes from docs/paper/results/*.json, produced by
% scripts/research_report.py from per-case rows; tables and figures are generated by
% scripts/make_figures.py. Sections marked RESULTS-PENDING are written only after the
% corresponding run exists - no number here is estimated.
\documentclass[11pt]{article}
\usepackage[margin=1in]{geometry}
\usepackage{booktabs,graphicx,amsmath,amssymb,hyperref,xcolor,enumitem,microtype}
\usepackage[numbers]{natbib}
\hypersetup{colorlinks=true,linkcolor=blue!50!black,citecolor=blue!50!black,urlcolor=blue!50!black}

\title{Deciding, Not Predicting: Symbolic Reasoning and Verification Between a Language
Model and Enterprise Policy Decisions}
\author{Praxis project\\\small (authors to be added)}
\date{}

\begin{document}
\maketitle

\begin{abstract}
Language models deployed as enterprise assistants are asked questions whose answers are
\emph{decisions}: whether an order can be returned, how much a customer is owed, whether
a remedy needs approval. A fluent model answers such questions by predicting a plausible
decision, and retrieval-augmented generation (RAG) does not change that: the model still
decides, now with policy text in its context. We study an alternative in which the
language model never decides. In \textsc{Praxis}, the model is a semantic parser that
maps a question into a knowledge pack's closed schema; authorization restricts which
records exist for the asker \emph{before} any reasoning; a three-valued rule engine
evaluates the policy in force on the case's own date over company records; and a
verification gate stands between the conclusion and the reply, withholding any answer
whose facts, provenance, calculation, policy version or wording do not check---a
guarantee which, we report, this evaluation never exercises. We introduce a benchmark of
procedurally generated
companies with counterfactual policy parameters, records placed on decision
boundaries, and gold answers from an independent reference implementation of the policy
text, validated by agreement with the rule engine and by mutation analysis. Across ten
question categories we compare the same open-weights model used alone, with retrieval,
with retrieval and self-verification, and with retrieval and a classical rule engine,
against \textsc{Praxis} using that model only as a parser, and we ablate each symbolic
mechanism. We report hallucination together with coverage, so that abstention is never
mistaken for reliability.
On 255 test questions, with the same open-weights model serving every system that uses
one, the architecture answers $0.945$ correctly at an unsupported-answer rate of $0.035$,
against $0.200$ and $0.302$ for that model with retrieval over the same corpus and $0.537$
and $0.247$ for it driving a classical rule engine. On a confirmatory split of five
further companies, frozen after the system was final and scored once, the configuration
whose parser was corrected using test-split failures falls to $0.835$ and $0.118$, and we
report that figure as the one that generalises; the configuration with no language model
in the loop replicates at $0.812$ and $0.133$, as does the retrieval baseline it is
measured against, at $0.216$ and $0.306$. Authorization replicates exactly: none of
fifteen restricted records disclosed, where moving authorization after reasoning discloses
thirteen, on both splits. Hallucination is reduced, not eliminated.
\end{abstract}

\section{Introduction}

An assistant for a support agent is asked: \emph{``My customer's laptop arrived damaged
twelve days ago and they want a replacement. What can I offer?''} The correct answer
depends on the order's delivery date, the product's stock, the returns policy
\emph{version} that governed the order on that date, the damage-reporting window that
version sets, the order total against an approval threshold, and the agent's
authorization to see the order at all. None of these is a matter of phrasing. A model that
answers well most of the time is not an acceptable decision procedure for this question,
because its errors are indistinguishable from its correct answers: both are fluent,
specific and confident.

Hallucination in this setting has a narrower and more testable meaning than in open-ended
generation: an answer is unsupported when it asserts a decision, an amount or a policy
value that the company's policies applied to the company's records do not produce, or when
it asserts anything definite where the knowledge supports nothing definite---a question
that omits a fact the policy needs, a topic no policy covers, a record the asker may not
see. This definition admits exact gold answers, and so admits measurement without human or
model judges.

We make three contributions.
\begin{enumerate}[leftmargin=*]
\item \textbf{An architecture} in which every decision is derived symbolically and the
language model's role is confined to mapping language into a closed, validated schema
and, optionally, to phrasing an already established conclusion under verification
(Section~\ref{sec:architecture}).
\item \textbf{A benchmark methodology} for enterprise decision questions that avoids the
three weaknesses of hand-written evaluation sets---a single dataset the system's author
knows, author-labelled gold answers, and policy values a model may know from pretraining
(Section~\ref{sec:benchmark}). Gold answers come from an independent reference
implementation; its agreement with the rule engine and its sensitivity to plausible
misreadings of each clause are measured, not assumed.
\item \textbf{A controlled comparison} holding the model, the knowledge and the questions
fixed while varying only the architecture, with coverage reported alongside
hallucination, paired significance tests, and component ablations
(Sections~\ref{sec:protocol}--\ref{sec:results}).
\end{enumerate}

Of the symbolic mechanisms the architecture contains, the ablations support two as doing
measurable work---authorization before reasoning, and three-valued logic together with
version ordering---and do not support two others. We claim the first pair and describe the
second as guarantees whose value this evaluation does not demonstrate, because a mechanism
that cannot be shown to matter should not be listed among a paper's contributions
(Sections~\ref{sec:results} and~\ref{sec:threats}).

We do not claim that the architecture eliminates hallucination. It moves the failure
surface: a decision can no longer be invented, but a question can still be misunderstood,
and the benchmark measures both.

\section{Related work}

\paragraph{Hallucination and grounding.} Surveys characterise hallucination as output
unsupported by the input or by world knowledge \citep{ji2023survey,huang2025survey}.
Retrieval augmentation \citep{lewis2020rag} supplies evidence but leaves the model to use
it; self-reflective retrieval \citep{asai2024selfrag} and chain-of-verification
\citep{dhuliawala2024cove} ask the model to check itself, and sampling-based detection
\citep{manakul2023selfcheckgpt} infers unreliability from inconsistency. Claim-level
factuality metrics decompose generations into atomic facts \citep{min2023factscore}. Our
setting differs in that the ground truth is a computation over structured records, which
permits exact scoring of decisions and amounts.

\paragraph{Offloading reasoning from the model.} Program-aided and program-of-thought
methods delegate arithmetic to an interpreter \citep{gao2023pal,chen2023pot}; neuro-symbolic
pipelines translate problems into logic for a solver \citep{olausson2023linc,pan2023logiclm}
or into programs over tables \citep{cheng2023binder}. These systems let the model write the
formal program. \textsc{Praxis} instead fixes the program---the company's rules---and lets the
model choose only among declared goals and input values, which removes the model's
ability to write a wrong rule while keeping its ability to understand paraphrase. The
broader neuro-symbolic agenda motivating this division of labour is surveyed by
\citet{garcez2023neurosymbolic} and \citet{kautz2022third}.

\paragraph{Questionable premises, time and abstention.} Questions with false assumptions
elicit confident wrong answers \citep{kim2023qa2,yu2023crepe}, and time-sensitive knowledge
defeats static parametric memory \citep{vu2023freshllms}. Policy \emph{versions} are a
sharper case: both versions are true documents, and only the case's date says which
applies. Selective prediction studies the trade-off between answering and abstaining
\citep{elyaniv2010selective,kamath2020selective}; we adopt its view that a system's
hallucination rate is uninterpretable without its coverage.

\section{Architecture}\label{sec:architecture}

A \emph{knowledge pack} holds a company's policy documents (clause-numbered, versioned), a
knowledge graph of records, and rules in a small logic language. Rules hold logic only;
every number---windows, fees, caps, thresholds---is a fact in the graph, and a policy
change is a new policy-version node rather than a rule edit. The pack declares
\emph{goals} (what may be decided), \emph{inputs} (what a person may state, with their
admissible values), \emph{access grants} (which record types each role may see, and
whether only their own), and \emph{verified inputs} (claims only a record can establish,
such as a supervisor's approval).

A question passes through the following stages; the language model participates only in
the first and, optionally, the last.

\begin{enumerate}[leftmargin=*]
\item \textbf{Understanding.} Either an embedding model matches the question to goal and
input descriptions (offline configuration), or the language model acts as a
\emph{semantic parser}: it returns one declared goal or \texttt{none}, a request kind, a
policy clause identifier for informational questions, and stated input values. Each
field is validated---an undeclared goal or enumeration value is discarded, a stated
number absent from the question is discarded---and the model never sees a record.
\item \textbf{Authorization before reasoning.} The knowledge graph is replaced by the
asker's authorized view, derived from the pack's access grants; every later stage,
including retrieval and verification, sees only that view.
\item \textbf{Routing.} The goal's rule structure selects policy reasoning, multi-hop
explanation, calculation, recommendation, retrieval of policy text, clarification, a
knowledge-gap report, or general conversation for questions outside the company's
knowledge.
\item \textbf{Grounding.} Record references are resolved and ownership checked; stated
facts that contradict a record are rejected and reported (the record is used); verified
inputs without a record are rejected; claims about policy values in the question are
checked against the parameters the proof reads.
\item \textbf{Three-valued reasoning.} A backward-chaining engine evaluates the goal with
\emph{unknown} distinct from \emph{false}. The returns-policy version is resolved on the
order's delivery date, not the current date. When a needed input is unknown, the reasoner
names it and the understanding stage is asked for exactly that input (a
neural-symbolic loop); if it remains unknown the answer is conditional and says what it
depends on.
\item \textbf{Verification.} Ten deterministic checks: fact validity, authorization,
policy applicability, policy version, symbolic provenance, re-derivation of calculations,
support for every stated number, contradictions, agreement of the wording with the
conclusion, and absence of unsupported claims. An answer that fails is withheld for
review, never released.
\item \textbf{Phrasing.} A deterministic realizer states the conclusion; optionally the
language model rephrases it, and the rephrasing is released only if it passes the same
checks, otherwise the deterministic text is used.
\end{enumerate}

\section{Benchmark}\label{sec:benchmark}

\subsection{Generated companies}
A seeded generator produces companies (\emph{worlds}) that share one schema and one set of
rule files but differ in every policy parameter and every record. Parameters are drawn
from ranges that include values no retailer publishes (for example a 23-day fashion
return window, a 7\% cancellation fee capped at 750), so answers drawn from pretraining
priors about typical retail policy are visibly wrong. Each world has two returns-policy
versions whose damage-reporting windows differ, with the change dated so that orders
delivered on \emph{both} sides of it fall between the two windows. Records are placed on
decision boundaries rather than sampled: orders exactly on, just inside and just outside
each return window; shipments exactly at and beyond the delay threshold; failed payments
with and without captured funds; high-value orders with and without a recorded
supervisor exception. Policy documents are rendered from the parameters, and the
superseded version remains in the corpus as a document index would keep it.

\subsection{Questions and gold answers}
Questions are instantiated from templates in ten categories: \emph{known}, \emph{paraphrased},
\emph{multi-hop} (records found through relations), \emph{calculation}, \emph{policy
version}, \emph{recommendation}, \emph{missing information}, \emph{false premise} (a false
policy value, a false record fact, or a claimed but unrecorded approval),
\emph{unauthorized} (a record belonging to someone else), and \emph{not covered} (plausible
company questions no policy answers). Canonical phrasings are shared; every other phrasing
exists in exactly one of the development and test splits, so repairs made on development
questions cannot fit test wording. Every template asks in its goal's polarity, so a yes or
no means the same thing for every system.

Gold answers are computed by a reference oracle implemented directly from the policy text
in ordinary code, independent of the rule language, the reasoner and the knowledge graph.
The oracle is three-valued: a question lacking a needed fact has a \emph{conditional}
gold answer naming the fact.

\subsection{Validity of the gold answers}
Two implementations written from one specification can share a misreading, and two
implementations can agree on records that never exercise the clause where they would
differ. We therefore report two checks. \emph{Agreement:} the rule engine and the oracle are
evaluated on every record of every world under every stated-reason variant used, bypassing
language understanding.
%% RESULTS: 3,914 of 3,914 evaluations agree (scripts/check_oracle.py; results/oracle_check.json).
They agree on all 3{,}914 evaluations. \emph{Mutation analysis:} twelve mutants of the
oracle each encode a plausible misreading of one clause---applying the newest (or oldest)
policy version to every order, ignoring the membership extension, making a window
exclusive, ignoring stock, charging a fee to exempt members, making the delay threshold
inclusive, crediting an unrecorded approval, ignoring the approval threshold, never
requiring receipts, ignoring service years, ignoring inspection. A mutant is \emph{killed}
if it disagrees with the rule engine on at least one record. All twelve are killed. The
first version of the generator left the newest-version mutant alive---no record fell
between the two versions' windows---which is how the version-boundary placement above
came to be; we report this because it shows the check detecting a real weakness.

\subsection{Splits and freezing}
The development split has 153 questions over three worlds; the test split has 255
questions over five disjoint worlds. Both were generated and committed before any system
was run on the test split; the scored file, not the generator, is the benchmark, and a test
asserts that regenerating it reproduces the committed questions and gold answers exactly.
Development questions were used to find and fix system faults; test questions were not
inspected for that purpose.

\subsection{A second knowledge pack}\label{sec:second-domain}
The architecture is schema-driven, and a claim of generality is worth what it can be
tested on. We therefore wrote a second pack for an unrelated domain---health
insurance---and asked it forty questions in the same ten categories, with gold answers
computed from the records by arithmetic performed outside the system.

The domain was chosen to exercise what retail does not: waiting periods that refuse a
claim for elapsed time rather than for the treatment; a deductible and a co-payment that
make a settlement a chain of arithmetic rather than one rule; exclusions no plan reaches;
a network status that changes the member's share; and a dependant who may see exactly one
other member's records.

Writing it located three assumptions in the implementation that the first pack could not
have revealed, each of which we report because their absence from the first domain is
precisely what a single-domain evaluation hides. The company node was identified by a
literal name, so every rule reading the current date silently found nothing and all date
arithmetic failed. Record ownership was resolved through hard-coded entity types and
predicates drawn from retail, so no insurance record reached its owner and every question
was refused as another person's. Three answer templates named the retail company outright.
All three are now declared by the pack; the original pack declares what the code had
assumed, and its results are unchanged.

\section{Experimental protocol}\label{sec:protocol}

\paragraph{Systems.} All model-based systems use the same open-weights instruction-tuned
model served locally with greedy decoding.
\begin{itemize}[leftmargin=*]
\item \textbf{LLM only}: the model with the asker's identity and the date.
\item \textbf{LLM + RAG}: retrieval of the top policy clauses (all versions) and of the
records the asker is authorized to see, by embedding similarity.
\item \textbf{LLM + RAG + self-check}: LLM + RAG followed by a second pass in which the model
verifies its draft against the same context---policy version, numbers, conditions,
authorization, missing facts, premises---and revises it.
\item \textbf{LLM + RAG + basic rules}: the model extracts goal, record and stated facts; a
classical engine (two-valued, newest policy version, statements accepted as given, no
verification) decides; the model phrases the result.
\item \textbf{Praxis (LLM parser)}: the architecture above with the model as semantic
parser only. \textbf{Praxis (LLM parser + phrasing)} additionally lets the model phrase
answers under verification.
\item \textbf{Praxis (offline encoder)}: embedding understanding, no language model.
\end{itemize}
Baselines return a JSON object with the answer text, an outcome (answered, needs more
information, which record, not covered, not permitted), and a decision, amount or remedy;
this structured output is what makes exact scoring possible and is requested identically of
every baseline.

\paragraph{Metrics.} \emph{Accuracy}: the outcome and value match gold. \emph{Hallucination
rate}: the answer is a wrong definite answer, a definite answer where gold is conditional,
not covered or not permitted, or discloses an unauthorized record. \emph{Coverage}: among
questions with a definite gold answer, the fraction answered definitively
(1 $-$ over-abstention). \emph{Selective accuracy}: accuracy among definite answers.
\emph{Unsupported numbers}: numbers in the answer text supported by no policy parameter,
relevant record, derived quantity or the question itself---identifiers, dates, names and
the integers 0--2 excluded---reported as a lower bound.

\paragraph{Statistics.} Rates carry 95\% percentile bootstrap intervals (2{,}000 resamples).
Systems are compared on paired per-question outcomes with exact McNemar tests
\citep{mcnemar1947}; ablation effects carry paired bootstrap intervals on the difference.

\paragraph{Ablations.} Each removes one mechanism everywhere it acts, through a switch read
at the point of use so that a later stage cannot silently repair the removal: the
verification gate; version resolution (newest version always); three-valued logic (unknown
as false); statement grounding (stated facts override records, claimed approvals count as
recorded); authorization before reasoning (every record visible).

\section{Results}\label{sec:results}
All numbers are the 255 test questions, one open-weights model serving every system that
uses one, and rates carry 95\% bootstrap intervals in Table~\ref{tab:main}. Paired
comparisons are exact McNemar tests on per-question outcomes.

\paragraph{The architecture, against language models.} \textsc{Praxis} with the model as
parser answers $0.945$ of the questions correctly and is unsupported on $0.035$. The same
model given retrieval over the same corpus reaches $0.200$ accuracy at $0.302$
hallucination; asked to check its own draft, $0.247$ and $0.282$; alone, $0.235$ and
$0.427$. Against retrieval the paired test is 200 questions to 10 on correctness
($p=5\times10^{-47}$) and 76 to 8 on hallucination ($p=5\times10^{-15}$). Against the same
model driving a classical two-valued rule engine---the strongest baseline, at $0.537$ and
$0.247$---it is 114 to 10 ($p=2\times10^{-23}$) and 60 to 6 ($p=3\times10^{-12}$).

These gaps are not bought with silence. \textsc{Praxis} answers $0.976$ of the questions
that have a definite answer, against $0.462$ for retrieval and $0.505$ for self-check, and
its selective accuracy is $1.000$: on this split, when it commits, it is right. The
baselines abstain roughly twice as often and are wrong more often when they do answer.

\paragraph{Understanding is where the remaining error lives.} The same architecture with
embedding understanding instead of a parser scores $0.796$ and $0.149$. Nothing else
differs: the same rules, the same authorization, the same gate. Replacing one component
moves hallucination by a factor of four ($p=4\times10^{-9}$, 29 questions to none), which
places the residual failure in mapping language onto the schema rather than in the
reasoning over it. The lexical-encoder configuration answers almost nothing and is
reported at $0.000$ coverage so that its $0.000$ hallucination is not mistaken for a
result.

\paragraph{Authorization.} Fifteen questions ask about a record the asker may not see. No
configuration that authorizes before reasoning discloses one. Moving authorization after
reasoning---every other mechanism intact---discloses $0.867$ of them, thirteen of fifteen,
and costs accuracy as well ($0.737$; $p=6\times10^{-5}$ on both). The two language-model
baselines that retrieve only what the asker may see leak nothing, as expected; extractive
retrieval, self-check and the rule-engine baseline each disclose $0.067$.

\paragraph{Ablations.} Removing statement grounding costs the most accuracy, $0.796$ to
$0.690$ ($p=1\times10^{-8}$), by accepting what a question asserts over what the records
say; it also lowers hallucination, to $0.094$, by refusing a quarter of the answerable
questions, which is why coverage is reported beside it. Two-valued logic costs $0.745$ and
raises hallucination to $0.188$ ($p=0.004$ both ways). Always applying the newest policy
version costs $0.761$ at $0.184$ ($p=0.004$).

Three ablations change nothing measurable. Kind agreement is exercised and does not reach
significance ($p=0.62$ on correctness, $p=0.25$ on hallucination; it is nominally
\emph{better} without, at $0.804$). The verification gate and the number check change no
answer at all, for the structural reason given in Section~\ref{sec:threats}: both police
prose the model wrote, and in this configuration the model writes none.

\paragraph{Cost.} The symbolic path answers in 78\,ms with embedding understanding and
380\,ms with the parser. The language-model baselines take 5.8 to 12.5 seconds, and the
most expensive of them---the model applying rules---is also the one that leaks. Unsupported
numbers appear in $0.012$ of \textsc{Praxis} answers against $0.153$ for the model alone
and $0.372$ for extractive retrieval.

\begin{table}[t]
\centering\small
\IfFileExists{figures/table_main.tex}{\resizebox{\linewidth}{!}{\begin{tabular}{lcccccc}
\toprule
System & Accuracy & Hallucination & Coverage & Selective acc. & Disclosure & Unsupp. numbers \\
\midrule
Praxis (offline encoder) · dev-iter0 & 0.804 \scriptsize[0.74, 0.86] & 0.105 \scriptsize[0.06, 0.16] & 0.889 \scriptsize[0.83, 0.94] & 0.938 \scriptsize[0.88, 0.97] & 0.000 \scriptsize[0.00, 0.00] & 0.105 \scriptsize[0.06, 0.16] \\
Extractive RAG · dev-iter0 & 0.000 \scriptsize[0.00, 0.00] & 0.118 \scriptsize[0.07, 0.17] & 0.000 \scriptsize[0.00, 0.00] & -- & 0.000 \scriptsize[0.00, 0.00] & 0.255 \scriptsize[0.18, 0.33] \\
Praxis (offline encoder) · dev-iter1 & 0.837 \scriptsize[0.77, 0.89] & 0.105 \scriptsize[0.06, 0.16] & 0.929 \scriptsize[0.88, 0.97] & 0.940 \scriptsize[0.90, 0.97] & 0.000 \scriptsize[0.00, 0.00] & 0.000 \scriptsize[0.00, 0.00] \\
\bottomrule
\end{tabular}
}}{}
\caption{Test split. Rates with 95\% bootstrap intervals.}
\label{tab:main}
\end{table}

\section{Confirmation on unseen companies}\label{sec:confirm}
The test split is 255 questions over five generated companies, and a fair objection is
that a system developed alongside a benchmark will score well on it. We therefore
generated five further companies, froze them after the test split had been scored and the
system was final, and scored them once. The result is 510 questions over ten disjoint
companies, with no question shared. What this adds is companies rather than phrasings:
the questions use the test split's wordings over different records, policy parameters and
boundary placements, so it measures generalisation to an unseen company and not to unseen
language.

\paragraph{What replicates.} The architecture with embedding understanding scores $0.812$
accuracy at $0.133$ hallucination, against $0.796$ and $0.149$ on the test split. The
authorization result replicates exactly: of fifteen questions about records the asker may
not see, the full system discloses none and the configuration that authorizes after
reasoning discloses thirteen, at $p=1\times10^{-4}$---the same thirteen-of-fifteen
pattern, on companies the system has never seen. That ablation also costs $0.056$ accuracy
here ($0.757$; 14 questions to none, $p=1\times10^{-4}$).

The comparison with a language model also replicates. Retrieval over the same corpus
scores $0.216$ at $0.306$, against $0.200$ and $0.302$ on the test split, and the
architecture beats it on these unseen companies by 168 questions to 16 on correctness
($p=4\times10^{-33}$) and 71 to 27 on hallucination ($p=1\times10^{-5}$).

\paragraph{What does not.} The configuration using the language model as parser scores
$0.835$ at $0.118$, against $0.945$ and $0.035$ on the test split. Its advantage over
embedding understanding falls from fifteen accuracy points to two, and with it the
significance: on the test split the parser is better on 40 questions to 2
($p=4\times10^{-10}$), and here on 7 to 1 ($p=0.07$), which we report as no measured
difference.

We report this as the principal finding of the confirmatory split, and the explanation is
methodological rather than mysterious. The parser's last defect was found by reading its
test-split failures: it read a question about the asker's own entitlement as a request for
the policy in general, and the correction was written after seeing which questions it got
wrong. The test-split figure is therefore not a generalisation estimate, and the
confirmatory figure is. We keep both, and treat $0.835$ as the number that describes the
parser configuration on a company it has not seen.

This is what a held-out confirmatory split is for, and it is worth stating that the
inflated figure would otherwise have been reported in good faith: the change that produced
it was measured in isolation, fixed ten questions and broke none, and carried $p=0.002$.
None of that makes it a generalisation result.

\begin{table}[t]
\centering\small
\IfFileExists{figures/table_confirm.tex}{\resizebox{\linewidth}{!}{\begin{tabular}{lcccccc}
\toprule
System & Accuracy & Hallucination & Coverage & Selective acc. & Disclosure & Unsupp. numbers \\
\midrule
LLM + RAG · confirm & 0.216 \scriptsize[0.17, 0.27] & 0.306 \scriptsize[0.25, 0.36] & 0.481 \scriptsize[0.41, 0.55] & 0.436 \scriptsize[0.34, 0.53] & 0.067 \scriptsize[0.00, 0.20] & 0.039 \scriptsize[0.02, 0.06] \\
Praxis (LLM parser) · confirm & 0.835 \scriptsize[0.79, 0.88] & 0.118 \scriptsize[0.08, 0.16] & 0.948 \scriptsize[0.91, 0.98] & 0.909 \scriptsize[0.87, 0.95] & 0.000 \scriptsize[0.00, 0.00] & 0.012 \scriptsize[0.00, 0.03] \\
Praxis (offline encoder) · confirm & 0.812 \scriptsize[0.76, 0.86] & 0.133 \scriptsize[0.09, 0.18] & 0.943 \scriptsize[0.91, 0.97] & 0.889 \scriptsize[0.84, 0.93] & 0.000 \scriptsize[0.00, 0.00] & 0.012 \scriptsize[0.00, 0.03] \\
\bottomrule
\end{tabular}
}}{}
\caption{Confirmatory split: five companies generated and frozen after the system was
final, scored once. Rates with 95\% bootstrap intervals.}
\label{tab:confirm}
\end{table}

\section{Calibration of the computed confidence}\label{sec:calibration}
The system never asks the model how sure it is. Confidence is computed in code from
retrieval relevance, rule coverage and the verification outcome, and no metric in this
paper reads it: it is recorded alongside correctness and analysed afterwards. That is a
claim about the design, and this section is what turns it into a measurement.

On the eighty-six question curated split, over ten equal-width bins, the expected
calibration error is $0.192$, the worst bin's gap $0.650$, and the Brier score $0.203$,
against an accuracy of $0.733$ and a mean confidence of $0.779$. Restricted to the
sixty-two questions the system answered rather than declined, the expected calibration
error is $0.179$ and the gap between mean confidence and accuracy widens to $0.156$: when
it answers, it says $0.914$ and is right $0.758$ of the time.

Two features of that distribution matter more than the summary statistic. Forty-two of
the eighty-six questions fall in the top bin, where the system claims $0.959$ and scores
$0.786$. And the scale is barely used: no question receives a confidence between $0.1$ and
$0.5$, so in practice the score takes about three values.

We therefore make the weaker claim the evidence supports. Computing confidence in code
rather than eliciting it from the model is defensible on its own terms---it is
reproducible, inspectable, and cannot be inflated by a fluent generator---but
\emph{computed} is not \emph{calibrated}, and nothing here licenses reading the number as
a probability of correctness. Fitting the weights on a development split and reporting on
a held-out one is the obvious remedy and is not done here.

\section{Error analysis}\label{sec:errors}
We read every failure on the second domain by hand: forty questions, thirty-four correct,
six wrong, and---the shape that matters before the detail---none of the six asserting
anything false. The hallucination rate on that split is $0.000$.

\emph{A stated fact was not recognised} (three of six). Each of these questions states, in
ordinary words, the one fact the records cannot hold, and the input detector scored the
paraphrase below threshold, so the reasoner declined to commit and answered conditionally.
This is a single defect seen three times: the input threshold is calibrated on the first
pack, and nothing has been tuned for the second. It is not a wrong answer. In each case
the system named the missing fact and made the answer contingent on it, which is what
three-valued logic is for; a system that had guessed would have been right and would have
been guessing.

\emph{The wrong goal was selected} (two of six). One question reached policy retrieval
instead of the settlement calculation. The other asserts a false premise---that an
excluded treatment is covered---and the system rejected the premise and stated that no
amount was due, by way of the wrong goal. The content is right; only the route is wrong,
and it is scored as a failure because a system that reaches the right words by the wrong
route will not keep doing so.

\emph{A question with no goal behind it} (one of six). The schema declares no goal for
whether a hospital is in the insurer's network, so the nearest declared goal won. This is
the same failure the first domain records as \emph{not covered}, and it is fixable in the
pack rather than in the architecture.

We leave one expectation deliberately unmet. A question using wording that sits between
two goals drew a clarification rather than a choice, which is defensible and possibly
better than guessing; the gold says otherwise, and was written before the run. Changing it
afterwards is how a benchmark stops meaning anything.

Across both domains the pattern is consistent: where this architecture fails, it fails
towards silence rather than towards invention.

\section{Threats to validity}\label{sec:threats}
\paragraph{Synthetic companies.} The generated worlds share one schema and one domain
(retail support and employee policy). A second pack in an unrelated domain
(Section~\ref{sec:second-domain}) shows the architecture carries, and located three
implementation assumptions that one domain had hidden; but that pack is forty questions
against the first domain's several hundred, its questions were written by the same author
in the same sitting as its schema, and it has no policy-version category. Its higher
score should not be read as evidence that the architecture generalises \emph{better}: the
two evaluations are not matched in size or difficulty, and only the first is generated
rather than written.
\paragraph{Confidence is not calibrated.} Confidence is computed rather than elicited, which
is a property of the design; Section~\ref{sec:calibration} measures it and finds it
overconfident and close to three-valued in practice. No result in this paper depends on
the number---no metric reads it---but a deployment that gated on it would be relying on a
quantity this evaluation shows to be poorly calibrated.
\paragraph{Mechanisms with no measured effect, and why.} Three of the six ablations change
nothing measurable, and the reason is not the same in each case, so we separate them
rather than reporting one null three times.

Removing the verification gate changes no answer. Removing the number check---which
refuses support to a claim whose figure does not appear in the evidence---changes no
answer either, and instrumenting the run shows why: in this configuration the claim
scorer is never called at all. Both belong to the same subsystem, and that subsystem
exists to police prose the model wrote. In the configuration evaluated here the model
writes no prose: answers are composed from the reasoning's conclusion, and the
claim check returns early on exactly that ground. The verification layer is therefore
insurance against a failure this configuration cannot have. That was our explanation of
the null, and it predicted that the gate would do work in the configuration where the
model phrases the answer under verification. We ran that configuration, with and without
the gate, over sixty questions: the model wrote the prose in 69 of 73 cases where its
wording could be compared with the deterministic one, and the gate changed no answer
there either.

So the explanation is wrong, or at least incomplete. On this data the verification gate
never withholds anything, including when there is model-written prose for it to withhold.
We report the prediction and its refutation rather than the explanation alone, because an
account of a null result that survives only until it is tested is not worth the reader's
confidence. What the gate is, on this evidence, is a guarantee that has never been
exercised; whether that is because the model's phrasing is tightly enough constrained by
the conclusion it is given, or because the checks are too permissive, this evaluation does
not distinguish.

Removing kind agreement is exercised and simply does not reach significance: three
questions are answered correctly without it against one with it, and three are
hallucinated without it against none with it, at $p=0.62$ and $p=0.25$. We report it as
no measured difference. A mechanism argued for on the grounds that it prevents a fluent
answer to a different question is not shown, on this data, to prevent one.
\paragraph{Template questions.} Questions are templated, and paraphrase diversity is bounded
by the templates. Disjoint development and test phrasings mitigate but do not remove this;
a study with questions written by practitioners is described in the pilot protocol
accompanying the code.
\paragraph{Shared specification.} Oracle and rules were written from the same policy text
by the same project; agreement shows consistency, not correctness of either reading.
Mutation analysis bounds the space of misreadings the benchmark can detect to those
enumerated.
\paragraph{Baseline strength.} Baselines use one open-weights model family at limited scale on
local hardware; stronger hosted models may hallucinate less. The comparison is controlled
for the model, which is its purpose, but absolute baseline rates should not be read as
properties of language models in general.
\paragraph{Structured baseline output.} Requiring JSON makes scoring exact but may cost small
models accuracy; answers whose JSON cannot be parsed are scored from what could be read, and
the parse-failure rate is reported.
\paragraph{Development on the same architecture.} The system's author designed the benchmark.
Freezing the test split before any run, and generating rather than writing its questions,
limit but do not eliminate this.

\section{Conclusion}
A language model that never decides can be held to a standard a fluent one cannot. With
the model confined to mapping a question onto a closed schema, every decision derived
symbolically, and authorization applied before any record enters a proof, the system
answers $0.945$ of 255 enterprise decision questions correctly at $0.035$ unsupported,
against $0.200$ and $0.302$ for the same model with retrieval, and discloses none of the
restricted records that the same architecture leaks thirteen of when authorization is
moved after reasoning. On five companies generated after the system was final, the
authorization result repeats unchanged and the headline accuracy does not: the
configuration whose parser was corrected against test-split failures scores $0.835$ there,
and that is the figure we would defend.

Two results temper that. Replacing only the component that reads the question moves
hallucination fourfold, so what remains of the error is in understanding rather than in
reasoning---and understanding is the one stage still performed by a model. And three of
six ablations change nothing measurable: the verification subsystem has no work to do in
a configuration where the model writes no prose, and kind agreement is not shown to
prevent the failure it was built to prevent. We report them because a mechanism that
cannot be shown to matter should not be claimed to.

The architecture reduces unsupported answers in enterprise question answering. It does
not eliminate them, and this paper does not claim a general result about hallucination.

\bibliographystyle{plainnat}
\bibliography{references}
\end{document}
